% Generated draft: requires mathematical and source review.
\documentclass[11pt]{article}
\usepackage[margin=25mm]{geometry}
\usepackage{fontspec}
\setmainfont{Latin Modern Roman}
\usepackage{amsmath,amssymb,mathtools,unicode-math}
\setmathfont{Latin Modern Math}
\usepackage{enumitem,needspace,xurl,hyperref}
\setlength{\parindent}{0pt}
\setlength{\parskip}{5pt}
\setlength{\emergencystretch}{3em}
\newcommand{\N}{\mathbb{N}}
\newcommand{\Z}{\mathbb{Z}}
\newcommand{\Q}{\mathbb{Q}}
\newcommand{\R}{\mathbb{R}}
\newcommand{\C}{\mathbb{C}}
\newcommand{\F}{\mathbb{F}}
\newcommand{\A}{\mathbb{A}}
\newcommand{\PP}{\mathbb{P}}
\newcommand{\E}{\mathbb{E}}
\newcommand{\eps}{\varepsilon}
\newcommand{\dd}{\,\mathrm{d}}
\begin{document}
\section*{UnsolvedMath ID: 700590}
\textbf{Source problem code:} catalogue entry
\subsection{Definitions and mathematical statement}
A formation of finite groups is understood here to be a nonempty class \(\mathfrak F\) of finite groups, closed under isomorphisms, homomorphic images, and finite subdirect products. Explicitly, the last condition means that if \(G_1,\ldots,G_n\in\mathfrak F\), where \(n\geq 1\), and \(H\leq G_1\times\cdots\times G_n\) projects surjectively onto every factor, then \(H\in\mathfrak F\). In particular, every formation contains the trivial group.

For a formation \(\mathfrak F\), let
\[
\operatorname{SubForm}(\mathfrak F)=\{\mathfrak A\subseteq\mathfrak F:\mathfrak A\text{ is a formation}\},
\]
ordered by inclusion. For any class \(\mathcal X\) of finite groups, write
\[
\operatorname{Form}(\mathcal X)=\bigcap\{\mathfrak H:\mathfrak H\text{ is a formation of finite groups and }\mathcal X\subseteq\mathfrak H\}.
\]
The meet and join of two subformations are
\[
\mathfrak A\wedge\mathfrak B=\mathfrak A\cap\mathfrak B,
\qquad
\mathfrak A\vee\mathfrak B=\operatorname{Form}(\mathfrak A\cup\mathfrak B).
\]
Because \(\mathfrak F\) contains \(\mathfrak A\cup\mathfrak B\), this generated formation is again a subformation of \(\mathfrak F\).

\textbf{Problem (A. N. Skiba, 1984; Kourovka Notebook 9.57).} Characterize all formations \(\mathfrak F\) of finite groups for which \(\operatorname{SubForm}(\mathfrak F)\) is a distributive lattice. Thus the required property is that, for every \(\mathfrak A,\mathfrak B,\mathfrak C\in\operatorname{SubForm}(\mathfrak F)\),
\[
\mathfrak A\wedge(\mathfrak B\vee\mathfrak C)
=(\mathfrak A\wedge\mathfrak B)\vee(\mathfrak A\wedge\mathfrak C),
\]
and, equivalently in a lattice, the dual distributive identity holds.

The supplied statement imposes no solubility, saturation, or other additional restriction on the formations. It does not define the term formation or specify whether the empty class is permitted; the formulation above uses the standard nonempty convention, which requires confirmation against the cited source.
\subsection{Short English statement}
Classify the formations of finite groups whose subformations form a distributive lattice, with intersection as meet and formation generation as join. The draft uses nonempty formations; the source's convention concerning the empty class remains to be checked.
\subsection{Sources}
\begin{itemize}
\item Local UnsolvedMath registry, permanent ID \(700590\), exact source problem code \texttt{catalogue entry}, title \emph{Kourovka Notebook Problem 9.57}; supplied upstream revision \texttt{d371ee917ac9647fbb06cd6fc529deec6f3dda82}.
\item \emph{The Kourovka Notebook (Unsolved Problems in Group Theory)}, original problem 9.57, attributed to A. N. Skiba, 9th issue (1984). The supplied record identifies the reference edition as the 21st issue edition (2026), arXiv:1401.0300v48.
\item The supplied collection description names E. I. Khukhro and V. D. Mazurov as editors of the Kourovka Notebook collection covering Issues 1--21 (1965--2026). These bibliographic details are supplied attribution and have not been independently verified.
\item The local registry labels the problem \emph{open}. This is a supplied status label, not an independently verified assessment of the literature.
\end{itemize}
\end{document}
